\chapter{Referencial Teórico}

Este capítulo apresenta os conceitos que fundamentam o trabalho sob a perspectiva de Sistemas de Informação. O foco recai sobre processamento inteligente de documentos, modelos multimodais, padronização de entrada e saída, validação de saídas estruturadas e avaliação experimental de artefatos computacionais.

\section{Processamento de Documentos Manuscritos}

O processamento de documentos manuscritos envolve a conversão de um artefato físico, sujeito a variações de escrita, captura, organização visual e marcação, em registros digitais que possam ser armazenados e utilizados por um sistema. Em provas objetivas manuscritas, a extração de informação exige reconhecer identificação do aluno, enunciados, numeração das questões, alternativas disponíveis, marcações de resposta, rasuras e ambiguidades visuais. A dificuldade aumenta quando a captura ocorre por fotografia, pois o sistema passa a lidar com inclinação, sombras, reflexos, baixa resolução, cortes parciais, caligrafias distintas e leiaute manuscrito irregular. O reconhecimento de manuscritos, tanto em cenários on-line quanto off-line, é um problema consolidado na área de análise de padrões e documentos \cite{plamondon2000handwriting,rassul2025handwritten}.

Soluções tradicionais baseadas em Reconhecimento Óptico de Marcas (OMR) tendem a funcionar melhor quando há um formulário padronizado e condições controladas de digitalização. Técnicas de OCR, como o Tesseract, e abordagens recentes de OMR, incluindo modelos leves e processamento com câmeras de smartphone, mostram diferentes caminhos para digitalização de documentos e marcações \cite{smith2007tesseract,mondal2023omrnet,hafeez2024omrsmartphone}. Já abordagens baseadas em entendimento de documentos podem interpretar simultaneamente texto, leiaute e imagem, reduzindo a dependência de templates rígidos \cite{xu2020layoutlm,kim2022donut}. Essa flexibilidade, entretanto, precisa ser avaliada empiricamente antes de ser incorporada a um sistema de informação, especialmente quando o documento é manuscrito.

\section{Modelos Multimodais}

Modelos multimodais de linguagem combinam processamento textual com interpretação visual, permitindo que uma imagem seja fornecida junto de uma instrução textual. No contexto deste trabalho, o modelo recebe a imagem da prova objetiva manuscrita e uma especificação do formato de saída esperado, retornando dados estruturados sobre questões identificadas, alternativas detectadas, respostas marcadas, identificação do aluno, rasuras, ambiguidades visuais e grau de confiança. Tarefas como VQA em documentos, pré-treinamento com texto e leiaute, e entendimento de documentos sem OCR indicam a relevância dos modelos multimodais para esse tipo de dado visual \cite{mathew2021docvqa,xu2020layoutlm,kim2022donut}.

Apesar do potencial técnico, a resposta de um modelo multimodal não deve ser tratada automaticamente como dado válido. Como esses modelos são probabilísticos, o sistema precisa considerar respostas incompletas, inconsistentes ou incompatíveis com o esquema esperado. Por isso, o fluxo proposto incorpora validação de JSON, registro de incertezas e revisão humana para casos que não atendem aos critérios de qualidade.

\section{Modelos Abertos e Modelos Proprietários}

A comparação entre modelos abertos e proprietários é relevante para decisões técnicas em Sistemas de Informação. Modelos abertos, frequentemente disponibilizados com pesos abertos e documentação pública, podem oferecer maior transparência sobre características gerais, possibilidade de execução local em cenários futuros e menor barreira para experimentação. Modelos proprietários, por outro lado, costumam ser disponibilizados como serviços fechados, com infraestrutura gerenciada pelo fornecedor, podendo oferecer estabilidade operacional, desempenho superior ou recursos adicionais.

Neste trabalho, ambos os modelos serão acessados por APIs hospedadas, de modo que a comparação não tem como foco hospedagem local, mas sim o comportamento dos modelos na tarefa de extração estruturada de informações em provas objetivas manuscritas. Como exemplo de enquadramento experimental, um modelo da família Gemma pode representar a categoria de modelo aberto ou \textit{open-weight}, enquanto um modelo da família Gemini pode representar a categoria proprietária, ambos consumidos por APIs documentadas pelo provedor \cite{google2026gemma,google2026gemmageminiapi,google2026geminiapi}. Essa distinção permite discutir graus de abertura, documentação, custo, latência, validade da saída estruturada e robustez em uma mesma tarefa \cite{solaiman2023gradient,bommasani2023transparency,ntia2024openweights,nist2023airmf}.

\section{Padronização de Entrada e Saída para Comparação de Modelos}

Para comparar modelos diferentes de forma controlada, é necessário padronizar as entradas, o prompt-base, o contrato de saída e as métricas coletadas. Neste trabalho, a camada de integração com os modelos não é tratada como contribuição principal, mas como mecanismo técnico para reduzir variações externas ao experimento. Assim, as imagens das provas, as instruções enviadas aos modelos e o esquema JSON esperado serão mantidos equivalentes para o modelo aberto e para o modelo proprietário.

A camada \textit{ModelAdapter} funciona como fronteira técnica entre o backend e as APIs dos modelos. Ela recebe uma requisição padronizada, chama o modelo configurado, interpreta a resposta bruta e retorna uma estrutura compatível com o esquema esperado pelo sistema. Essa separação apoia o protocolo experimental, pois permite registrar métricas de execução e manter a mesma rotina de validação humana para ambos os modelos.

\section{Saída Estruturada e Validação Humana}

Em um sistema transacional, uma resposta textual livre produzida por IA não é suficiente. O resultado precisa ser convertido em uma estrutura previsível, validada contra regras formais e persistida com rastreabilidade. Por isso, o trabalho adota uma saída em JSON contendo, no mínimo, identificador da submissão, respostas extraídas, campos de confiança, erros de validação e indicação de revisão humana.

A validação humana é tratada como parte do fluxo, e não como uma etapa externa. Quando o modelo retorna baixa confiança, resposta ausente, alternativa ambígua ou JSON inválido, o sistema apresenta o caso em uma tela de validação. A correção feita pelo usuário é armazenada junto do resultado original, permitindo comparar a saída do modelo com o resultado final validado. Essa abordagem se alinha a recomendações de interação humano-IA, auditoria algorítmica e gestão de riscos em sistemas baseados em IA \cite{amershi2019humanai,raji2020audit,nist2023airmf}.

\section{Design Science Research e Avaliação de Artefatos}

A metodologia Design Science Research (DSR) é adequada para este trabalho por orientar a construção e avaliação de artefatos voltados à solução de problemas em Sistemas de Informação \cite{sposito2024design}. O artefato proposto não é apenas uma interface de usuário, mas um fluxo operacional composto por aplicativo, backend, camada de integração com modelos, banco de dados e mecanismos de validação.

A avaliação do artefato será conduzida de forma experimental e comparativa, observando métricas técnicas e operacionais. O Framework for Evaluation in Design Science (FEDS) apoia a definição de avaliações formativas e somativas ao longo do ciclo de desenvolvimento \cite{feds_framework2026}. Neste trabalho, a avaliação final deverá indicar qual abordagem técnica se mostrou mais adequada para o problema definido, considerando não apenas acurácia, mas também custo, latência, taxa de saída válida e necessidade de intervenção humana.

\section{Trabalhos Relacionados}

Diversos trabalhos abordam a correção automática de avaliações objetivas por meio de Reconhecimento Óptico de Marcas (OMR), processamento de imagem e captura por câmera. O Eyegrade, proposto por Fisteus, Pardo e García, utilizou técnicas de visão computacional de baixo custo para corrigir exames de múltipla escolha com webcam \cite{fisteus2013eyegrade}. Zampirolli, Quilici-Gonzalez e Neves também investigaram a correção automática de testes de múltipla escolha com câmeras digitais \cite{zampirolli2013automatic}, incluindo trabalhos posteriores em dispositivos Android \cite{zampirolli2015android}. Abordagens recentes, como OMRNet e métodos tolerantes a falhas com câmeras de smartphone, reforçam a relevância desse campo \cite{mondal2023omrnet,hafeez2024omrsmartphone}. Em geral, porém, essas soluções assumem formulários ou folhas de resposta mais estruturadas.

No contexto brasileiro, o sistema Minha Prova propôs a automatização do processo avaliativo escolar \cite{silva2019minhaprova}, enquanto o trabalho Avalia investigou a correção automatizada de atividades escolares utilizando OCR e inteligência artificial \cite{oliveira2024avalia}. Mais recentemente, Almeida et al. apresentaram um sistema para geração e correção automática de provas utilizando IA generativa, arquitetura de microsserviços, modelos locais via Ollama, banco vetorial e RAG \cite{almeida2025geracao}. Esses trabalhos indicam o crescimento do uso de IA em processos avaliativos, mas diferem do presente estudo por concentrarem-se na aplicação educacional, na geração de questões ou na correção automatizada em sentido amplo.

Pesquisas em \textit{document understanding} e modelos multimodais, como LayoutLM, DocVQA, Donut e DocLLM, demonstram a evolução de técnicas capazes de combinar texto, imagem e leiaute \cite{xu2020layoutlm,mathew2021docvqa,kim2022donut,docllm2024}. Estudos sobre reconhecimento de manuscritos com LLMs multimodais também indicam o uso recente desses modelos em documentos visualmente complexos \cite{li2024handwritingmllm}. Esses trabalhos fornecem base técnica para a extração estruturada de documentos, mas não tratam especificamente da comparação experimental entre modelos multimodais abertos e proprietários, ambos acessados por APIs hospedadas, em uma tarefa de extração estruturada de provas objetivas manuscritas.

Dessa forma, a diferença central deste TCC está na avaliação comparativa de modelos multimodais sob um fluxo experimental padronizado, utilizando o mesmo conjunto de imagens, o mesmo prompt-base, o mesmo contrato JSON e as mesmas métricas. A contribuição não se limita ao desenvolvimento de uma aplicação de correção, mas à análise técnica da validade da saída estruturada, da necessidade de validação humana e dos trade-offs de acurácia, custo, latência e robustez.

\section{Síntese do Referencial}

O referencial teórico sustenta que a contribuição deste trabalho está na avaliação comparativa de modelos multimodais em um sistema de informação, com ênfase em extração estruturada de documento manuscrito, padronização de entrada e saída, contrato JSON, validação humana e mensuração experimental. A aplicação em provas objetivas manuscritas oferece um caso concreto para medir o comportamento dos modelos em documentos físicos e visualmente heterogêneos.
